\chapter{Complexity — Adder Language}
%%% generated by the blueprint pipeline from TCSlib.Complexity.CircuitComplexity.AdderLanguage
\section{Overview}

Declarations of the Lean module \texttt{TCSlib.Complexity.CircuitComplexity.AdderLanguage}.

\section{Declarations}

\begin{definition}[Finite list of binary digits]
\lean{BoolCircuit.natToBits}
\label{BoolCircuit.natToBits}
\leanok
For natural numbers $k$ and $m$, form the length-$k$ list whose entry at index $j$, for
$0\le j<k$, is the $j$th binary digit of $m$ (with the least significant digit at index
zero).
\end{definition}

\begin{theorem}[Length of a fixed-width bit list]
\lean{BoolCircuit.length_natToBits}
\label{BoolCircuit.length_natToBits}
\leanok
\uses{BoolCircuit.natToBits}
\difficulty{1}
For any natural numbers $k$ and $m$, the list of the lowest $k$ binary digits of $m$ has
length $k$.
\end{theorem}

\begin{theorem}[Value of the low-order bits]
\lean{BoolCircuit.bitsVal_natToBits}
\label{BoolCircuit.bitsVal_natToBits}
\leanok
\uses{BoolCircuit.bitsVal, BoolCircuit.natToBits}
\difficulty{5}
Let $k,m$ be natural numbers with $m<2^k$. The value of the length-$k$ Boolean list
consisting of the binary digits of $m$ from positions $0$ through $k-1$ is $m$.
\end{theorem}

\begin{definition}[Adder circuit family]
\lean{BoolCircuit.adderFamily}
\label{BoolCircuit.adderFamily}
\leanok
\uses{BoolCircuit.DAGCircuitFamily, BoolCircuit.adderCircuit, BoolCircuit.constCircuit}
For each natural number $n$, assign the adder circuit for inputs of length $n$ when $n$
is divisible by $3$, and the constant-false circuit for inputs of length $n$ otherwise.
\end{definition}

\begin{definition}[Language of binary addition]
\lean{Language.addition}
\label{Language.addition}
\leanok
\uses{BoolCircuit.bitsVal}
A Boolean word belongs to this language exactly when it is the concatenation of three
Boolean lists $a,b,c$ of equal length and the value represented by $c$ is the sum of the
values represented by $a$ and $b$.
\end{definition}

\begin{theorem}[Adder family fan-in bound]
\lean{BoolCircuit.adderFamily_hasFaninTwo}
\label{BoolCircuit.adderFamily_hasFaninTwo}
\leanok
\uses{BoolCircuit.DAGCircuitFamily.HasFaninTwo, BoolCircuit.adderCircuit_isFaninTwo, BoolCircuit.adderFamily, BoolCircuit.constCircuit_isFaninTwo}
\difficulty{4}
Every circuit in the adder family is well formed: each gate has distinct input vertices,
and every negation gate has exactly one input. Moreover, each gate has at most two input
vertices.
\end{theorem}

\begin{theorem}[Adder family size bound]
\lean{BoolCircuit.adderFamily_size_le}
\label{BoolCircuit.adderFamily_size_le}
\leanok
\uses{BoolCircuit.DAGCircuit.size, BoolCircuit.adderCircuit_size, BoolCircuit.adderFamily, BoolCircuit.constCircuit_size}
\difficulty{4}
For every natural number $n$, the size of the adder family circuit on inputs of length
$n$ is at most $6n+4$.
\end{theorem}

\begin{theorem}[Language recognized by the adder family]
\lean{BoolCircuit.adderFamily_language}
\label{BoolCircuit.adderFamily_language}
\leanok
\uses{BoolCircuit.DAGCircuitFamily.language, BoolCircuit.DAGCircuitFamily.mem_language_iff, BoolCircuit.adderCircuit_eval_append, BoolCircuit.adderFamily, BoolCircuit.constCircuit_eval, Language.addition}
\difficulty{5}
The language recognized by the adder family consists exactly of the Boolean strings that
encode triples of bit strings $a,b,c$ with $a$ and $b$ having the same length as $c$,
where the binary value of $c$ is the sum of the binary values of $a$ and $b$.
\end{theorem}

\begin{theorem}[Linear size for addition]
\lean{Language.addition_inSIZE}
\label{Language.addition_inSIZE}
\leanok
\uses{BoolCircuit.adderFamily, BoolCircuit.adderFamily_hasFaninTwo, BoolCircuit.adderFamily_language, BoolCircuit.adderFamily_size_le, Language.InSIZE, Language.addition}
\difficulty{1}
The language of Boolean-word encodings of addition is recognized by a family of
well-formed acyclic Boolean circuits with fan-in at most two such that, on inputs of
length $n$, the circuit size is at most $6n+4$.
\end{theorem}

\begin{theorem}[Polynomial size for addition]
\lean{Language.addition_inPPoly}
\label{Language.addition_inPPoly}
\leanok
\uses{Language.InPPoly, Language.InSIZE.inPPoly, Language.addition, Language.addition_inSIZE}
\difficulty{3}
The language of Boolean-word encodings of addition is recognized by a family of
well-formed acyclic Boolean circuits with fan-in at most two whose size is bounded by a
polynomial in the input length.
\end{theorem}

\begin{theorem}[Fixed-width binary addition words]
\lean{Language.natToBits_mem_addition}
\label{Language.natToBits_mem_addition}
\leanok
\uses{BoolCircuit.bitsVal_natToBits, BoolCircuit.natToBits, Language.addition}
\difficulty{3}
Let $m,n,k$ be natural numbers such that $m+n<2^k$. The concatenation of the $k$-bit
representations of $m$, $n$, and $m+n$ belongs to the Boolean addition language, whose
words encode three equally long bit strings with the value of the third equal to the sum
of the first two.
\end{theorem}

\begin{theorem}[Fixed-width binary representations of addition]
\lean{Language.exists_natToBits_mem_addition}
\label{Language.exists_natToBits_mem_addition}
\leanok
\uses{BoolCircuit.natToBits, Language.addition, Language.natToBits_mem_addition}
\difficulty{3}
For any natural numbers $m$ and $n$, there is a width $k$ such that the concatenation of
their length-$k$ binary representations and the length-$k$ binary representation of
$m+n$ belongs to the language of addition; in this language, the three blocks have equal
length and the value represented by the third block is the sum of the values represented
by the first two.
\end{theorem}
